\documentclass[11pt,a4paper]{letter}
\usepackage[utf8]{inputenc}
\usepackage[T1]{fontenc}
\usepackage[margin=2.5cm]{geometry}
\usepackage{hyperref}

\signature{Luiz Diego Vidal Santos\\
Programa de P\'{o}s-Gradua\c{c}\~{a}o em Planejamento Territorial\\
Universidade Estadual de Feira de Santana\\
Feira de Santana, BA 44036-900, Brazil\\
\texttt{ldvsantos@uefs.br}\\
ORCID: 0000-0001-8659-8557}

\date{\today}

\begin{document}

\begin{letter}{The Editor\\
\textit{Environmental Earth Sciences}\\
Springer Nature}

\opening{Dear Editor,}

We are pleased to submit the enclosed manuscript entitled \textbf{``A multiplicative correction of USLE erodibility for tropical Plinthosols: parametric identifiability and physical consistency''} for consideration as an Original Research Article in \textit{Environmental Earth Sciences}.

The erodibility factor ($K$) of the Universal Soil Loss Equation remains the most widely used metric of soil susceptibility to water erosion, yet its estimation via the standard nomograph is restricted to surface-layer attributes and systematically underestimates erosive susceptibility in tropical soils governed by subsurface processes. Plinthosols, which are widespread across tropical landscapes, combine subsurface hydraulic impedance, aluminum phytotoxicity and geotechnical vulnerability --- three convergent mechanisms absent from the original $K$-factor formulation.

This study proposes a multiplicative correction factor ($K_{\mathrm{plint}}$) that translates these three mechanisms as independent amplifiers of $K_{\mathrm{RUSLE}}$, with five free parameters. Through five computational tests anchored in field-measured properties of a Haplic Plinthosol from the Coastal Tablelands of Sergipe (northeastern Brazil), complemented by virtual profiles of a Ferralsol and an Acrisol, we demonstrate that the three hydraulic-biological parameters are recoverable with bias below 15\% even under 30\% observational noise, while Sobol global sensitivity analysis confirms that subsurface hydraulic impedance governs 63--78\% of amplification variance in the Plinthosol. The Ferralsol negative control consistently yields amplification near unity, validating the multi-class experimental design, and Green--Ampt infiltration modeling coupled with simplified WEPP-type erosion simulation preserves the pedological hierarchy across soil classes.

To our knowledge, no previous study has simultaneously incorporated subsurface hydraulic impedance, aluminum phytotoxicity and geotechnical vulnerability into a correction factor for $K$, nor evaluated the parametric identifiability of such a formulation prior to field calibration. The proposed approach may contribute to bridging the gap between the operational simplicity of USLE and the mechanistic complexity of tropical erosion processes.

We believe this manuscript aligns with the scope of \textit{Environmental Earth Sciences}, particularly its coverage of water--soil interaction, land degradation processes, soil erosion modeling and geotechnical aspects of earth-surface dynamics. The simulation scripts, input parameters and generated figures are publicly available at Zenodo (\url{https://doi.org/10.5281/zenodo.19483403}).

This manuscript has not been published previously, is not under consideration for publication elsewhere, and all authors have approved the submitted version. There are no competing interests to declare.

\closing{Sincerely,}

\end{letter}
\end{document}
